\section{制約付き最適化}

ここまでで、線形系、非線形モデル、推定、最適フィードバックがそろった。最後に、現実の制約を正面から扱う。入力には上限があり、状態にも越えてはいけない範囲がある。制約付き最適化の標準形は
\begin{align}
  \min_z\quad &\frac{1}{2}z^\T H z+g^\T z,\\
  \text{subject to}\quad &Az=b,\\
  &Cz\le d
\end{align}
である。$H>0$ なら凸二次計画問題であり、局所最適解が大域最適解になる。

等式制約だけなら、ラグランジュ乗数 $\lambda$ を使って
\begin{equation}
  \mathcal{L}(z,\lambda)=\frac{1}{2}z^\T Hz+g^\T z+\lambda^\T(Az-b)
\end{equation}
を作る。最適性条件は
\begin{equation}
  \begin{bmatrix}H&A^\T\\ A&0\end{bmatrix}
  \begin{bmatrix}z\\ \lambda\end{bmatrix}
  =
  \begin{bmatrix}-g\\ b\end{bmatrix}
\end{equation}
である。不等式制約が入ると、KKT 条件には相補性
\begin{equation}
  \mu_i(Cz-d)_i=0,
  \qquad
  \mu_i\ge0
\end{equation}
が加わる。制約が効いている成分だけに乗数が現れる。

\section{MPC}

MPC はここで初めて自然に現れる。離散時間モデル、有限ホライズン最適制御、状態推定、入力制約、数値最適化をすでに見たので、MPC は突然の新概念ではなく、それらを毎サンプル結び直す方法として読める。

離散時間線形系
\begin{equation}
  x_{k+1}=A_dx_k+B_du_k
\end{equation}
に対し、ホライズン $N$ の MPC は
\begin{align}
  \min_{x_i,u_i}\quad &
  \sum_{i=0}^{N-1}(x_i^\T Qx_i+u_i^\T Ru_i)+x_N^\T P_fx_N,\\
  \text{subject to}\quad &x_{i+1}=A_dx_i+B_du_i,\\
  &x_i\in\mathcal{X},\quad u_i\in\mathcal{U},\\
  &x_0=x_k
\end{align}
を解く。得られた入力列の最初の $u_0$ だけを実機へ加え、次のサンプルで状態を測り直して問題を解き直す。これを後退ホライズン制御という。

終端コスト $P_f$ に LQR のリカッチ解を使うと、ホライズンの先に続くコストを近似できる。終端制約
\begin{equation}
  x_N\in\mathcal{X}_f
\end{equation}
を置けば、制約を守りながら安定性を示しやすくなる。MPC の難しさは、最適性だけでなく、毎回解が存在する実行可能性を保つ点にある。

\section{NMPC と数値解法}

非線形モデルをそのまま予測に使うと、非線形モデル予測制御になる。
\begin{align}
  \min_{x(\cdot),u(\cdot)}\quad &
  \int_0^T \ell(x(t),u(t))\,\dd t+V_f(x(T)),\\
  \text{subject to}\quad &\dot{x}=f(x,u),\\
  &g(x,u)\le0
\end{align}
を考える。これを有限次元問題へ変える代表的な方法が多重射撃法である。区間ごとに状態 $x_i$ と入力 $u_i$ を変数として持ち、数値積分で得た写像 $\Phi$ に対して
\begin{equation}
  x_{i+1}-\Phi(x_i,u_i)=0
\end{equation}
を制約として課す。

SQP は、非線形最適化問題を現在の軌道のまわりで二次計画問題に近似し、解いた方向へ更新する方法である。コストの二階微分を厳密に使う代わりに、残差 $r(z)$ に対する最小二乗構造
\begin{equation}
  \frac{1}{2}\norm{r(z)}^2
\end{equation}
から
\begin{equation}
  \nabla^2 J(z)\simeq J_r(z)^\T J_r(z)
\end{equation}
と近似するのがガウス・ニュートン法である。

実時間反復法では、各サンプルで SQP を完全収束させず、一回の線形化と一回の QP 解法だけで入力を更新する。前サンプルの解を初期値として持ち越せば、短いサンプリング周期でも実用的な更新ができる。自動微分は、$f$ や $\ell$ のヤコビアンを手計算せず、計算グラフから機械的に正確な微分を得るための技術である。

\section{制御配分}

多入力系では、制御器が欲しい一般化力 $w$ を出し、それを実際のアクチュエータ入力 $u$ へ変換する問題が生じる。線形近似で
\begin{equation}
  w=Bu
\end{equation}
と書けるなら、$B\in\R^{p\times m}$ が配分行列である。$m>p$ の冗長系では、同じ $w$ を作る $u$ が無数にある。

最小ノルム解はムーア・ペンローズ擬似逆行列で
\begin{equation}
  u=B^+w,
  \qquad
  B^+=B^\T(BB^\T)^{-1}
\end{equation}
と書ける。ただし、これは入力制限を考えない。一般解は
\begin{equation}
  u=B^+w+(I-B^+B)z
\end{equation}
であり、第二項は合力・合モーメントを変えないヌル空間入力である。ここに消費電力低減や余裕確保の目的を入れられる。

制約を入れるなら
\begin{align}
  \min_u\quad &\frac{1}{2}\norm{Wu}^2,\\
  \text{subject to}\quad &Bu=w,\\
  &u_{min}\le u\le u_{max}
\end{align}
を解く。アクチュエータが飽和したとき、単純な擬似逆行列では一部の入力が上限を越える。制約付き配分は計算量が増えるが、実機の安全性と性能を同時に扱える。

\section{サンプリング、遅れ、飽和}

制御器は連続時間で考えても、実装は離散時間で動く。センサ取得、推定、最適化、通信、アクチュエータ応答には遅れがある。入力遅れを一サンプルとみなすなら
\begin{equation}
  x_{k+1}=A_dx_k+B_du_{k-1}
\end{equation}
であり、状態に過去入力を加えて拡大系
\begin{equation}
  \begin{bmatrix}x_{k+1}\\ u_k\end{bmatrix}
  =
  \begin{bmatrix}A_d&B_d\\0&0\end{bmatrix}
  \begin{bmatrix}x_k\\ u_{k-1}\end{bmatrix}
  +
  \begin{bmatrix}0\\ I\end{bmatrix}u_k
\end{equation}
として扱える。

量子化は連続値を有限の刻みに丸める。刻み幅を $\Delta$ とすると、丸め誤差は理想化すれば
\begin{equation}
  \abs{e_q}\le\frac{\Delta}{2}
\end{equation}
である。高ゲイン制御では小さな量子化ノイズが入力を揺らすことがある。センサノイズ、計算遅れ、飽和、量子化は、最後に付け足す実装上の注意ではなく、帯域やゲインを決める制約である。

飽和を含む入力は
\begin{equation}
  u=\operatorname{sat}(u_c)
  =\min\{u_{max},\max(u_{min},u_c)\}
\end{equation}
と書ける。飽和中は、線形制御で仮定した入力が実際には入っていない。積分器、オブザーバ、MPC の状態予測は、飽和後の実入力を使って更新する必要がある。

\section{階層的な実時間制御}

複雑な機械を一つの巨大な制御器だけで動かすのは難しい。時間スケールで分けると設計しやすくなる。遅い層では目標軌道や制約付き最適化を扱い、中間層では状態フィードバックや配分を行い、速い層では電流、速度、力などの局所ループを閉じる。

たとえば上位の MPC が一般化力 $w_k$ を決め、中位の配分がアクチュエータ入力 $u_k$ へ変換し、下位の PI ループが実際のモータ速度を追従させる。推定器は各層へ状態を供給し、ログと診断はモデル誤差を見つける。全体は
\begin{equation}
  \text{推定}\rightarrow \text{予測}\rightarrow \text{配分}\rightarrow \text{局所制御}\rightarrow \text{対象}
\end{equation}
という循環になる。

階層化しても、各層は独立ではない。上位が速すぎる目標を出せば下位は飽和し、下位の遅れを無視すれば上位の予測は外れる。よい実装では、上位の制約に下位の限界を入れ、下位の状態を上位へ戻す。制御工学の各道具は、単独で完結する章ではなく、この循環の中で互いに役割を持つ。

\section{制約付き最適化の条件を導く}

\begin{theorem}[KKT 条件]
制約付き最適化問題
\begin{align}
  \min_z\quad &J(z),\\
  \text{subject to}\quad &h(z)=0,\\
  &g(z)\le0
\end{align}
が適当な正則性条件を満たすとき、局所最適解 $z^*$ には乗数 $\lambda$、$\mu\ge0$ が存在して
\begin{align}
  \nabla J(z^*)+\nabla h(z^*)^\T\lambda+\nabla g(z^*)^\T\mu&=0,\\
  h(z^*)&=0,\\
  g(z^*)&\le0,\\
  \mu_i g_i(z^*)&=0
\end{align}
を満たす。これを Karush--Kuhn--Tucker 条件という。
\end{theorem}

\begin{derivation}
等式制約だけなら、制約面の上で $J$ が最小になる点では、許された微小変化方向に沿った $J$ の一次変化が 0 になる。これは $\nabla J$ が制約面の接空間に直交することを意味する。制約 $h(z)=0$ の法線方向は $\nabla h(z)$ で張られるので
\begin{equation}
  \nabla J(z)+\nabla h(z)^\T\lambda=0
\end{equation}
と書ける。不等式制約では、効いていない制約は局所的に無視できるので乗数が 0 になる。効いている制約では $g_i(z)=0$ で乗数が現れる。この二つを一つの式で表すのが
\begin{equation}
  \mu_i g_i(z)=0,
  \qquad
  \mu_i\ge0
\end{equation}
である。
\end{derivation}

\section{MPC を行列の二次計画問題へ落とす}

線形 MPC では、状態列を入力列で表すと二次計画問題になる。簡単のため $x_{k+1}=Ax_k+Bu_k$、ホライズン $N=2$ とする。
\begin{align}
  x_1&=Ax_0+Bu_0,\\
  x_2&=Ax_1+Bu_1=A(Ax_0+Bu_0)+Bu_1\\
      &=A^2x_0+ABu_0+Bu_1.
\end{align}
入力列を
\begin{equation}
  U=\begin{bmatrix}u_0\\u_1\end{bmatrix}
\end{equation}
とおくと、状態列は
\begin{equation}
  \begin{bmatrix}x_1\\x_2\end{bmatrix}
  =\begin{bmatrix}A\\A^2\end{bmatrix}x_0+
  \begin{bmatrix}B&0\\AB&B\end{bmatrix}U
\end{equation}
と書ける。評価関数は $U$ の二次式
\begin{equation}
  J(U)=\frac{1}{2}U^\T HU+g^\T U+\text{定数}
\end{equation}
になる。入力制約や状態制約も、同じ代入によって
\begin{equation}
  C_UU\le d_U
\end{equation}
の形へ集約できる。したがって線形 MPC は、毎サンプル「現在の $x_0$ をパラメータに持つ QP」を解いている。

\begin{derivation}
擬似逆行列の最小ノルム解も省略せずに見る。$Bu=w$ を満たす中で $\frac{1}{2}u^Tu$ を最小にする。ラグランジアンを
\begin{equation}
  \mathcal{L}(u,\lambda)=\frac{1}{2}u^Tu+\lambda^T(Bu-w)
\end{equation}
と置く。停留条件は
\begin{align}
  \frac{\pp\mathcal{L}}{\pp u}&=u+B^T\lambda=0,\\
  Bu-w&=0.
\end{align}
一つ目より $u=-B^T\lambda$。これを制約へ代入して
\begin{align}
  B(-B^T\lambda)&=w,\\
  -BB^T\lambda&=w,\\
  \lambda&=-(BB^T)^{-1}w.
\end{align}
したがって
\begin{equation}
  u=B^T(BB^T)^{-1}w=B^+w
\end{equation}
である。ここで $B$ が行フルランクであることを使った。
\end{derivation}

\section{実装条件をモデルへ入れる}

入力飽和を単なる注意として扱うと、理論と実機が離れる。計算入力 $u_c$ と実入力 $u$ の関係を
\begin{equation}
  u=\operatorname{sat}(u_c)
\end{equation}
と明示すれば、推定器や予測器は $u_c$ ではなく実際に対象へ入った $u$ で更新すべきだと分かる。遅れも同じである。一サンプル遅れがあれば、対象が見る入力は $u_{k-1}$ なので
\begin{equation}
  x_{k+1}=A_dx_k+B_du_{k-1}
\end{equation}
になる。過去入力を状態へ含めると
\begin{equation}
  \xi_k=\begin{bmatrix}x_k\\u_{k-1}\end{bmatrix}
\end{equation}
として
\begin{equation}
  \xi_{k+1}=\begin{bmatrix}A_d&B_d\\0&0\end{bmatrix}\xi_k+
  \begin{bmatrix}0\\I\end{bmatrix}u_k
\end{equation}
を得る。これは経験的な補正ではなく、遅れを含む対象を別の状態空間モデルへ書き直した結果である。
